% TODO make hw stylesheet
\documentclass[10pt]{scrartcl}
\usepackage[a4paper, total={6in, 8in}]{geometry}
\usepackage[usenames,dvipsnames,svgnames]{xcolor}
\usepackage[shortlabels]{enumitem}
\usepackage[framemethod=TikZ]{mdframed}
\usepackage{amsmath,amssymb,amsthm}
\usepackage[colorlinks]{hyperref}
\usepackage{multicol}
\usepackage{tabularx}

\hypersetup{
	colorlinks=true,
	linkcolor=blue,
	filecolor=magenta,      
	urlcolor=magenta,
	pdftitle={MAT 322 HW}
}

\usepackage{microtype}
\usepackage{mathtools}
\usepackage[headsepline]{scrlayer-scrpage}
\usepackage{thmtools}
\usepackage[textsize=footnotesize]{todonotes}

\mdfdefinestyle{mdbluebox}{roundcorner=10pt,innerbottommargin=9pt,
	linecolor=blue,backgroundcolor=TealBlue!5,}
\declaretheoremstyle[headfont=\sffamily\bfseries\color{MidnightBlue},
mdframed={style=mdbluebox},]{thmbluebox}

\mdfdefinestyle{mdbloobox}{frametitlefont=\bfseries,innerbottommargin=8pt,
	nobreak=true,backgroundcolor=TealBlue!5,linecolor=blue,}
\declaretheoremstyle[headfont=\bfseries\color{MidnightBlue},
mdframed={style=mdbloobox},headpunct={\\[3pt]},postheadspace=0pt,]{thmbloobox}

\mdfdefinestyle{mdredbox}{frametitlefont=\bfseries,innerbottommargin=8pt,
	nobreak=true,backgroundcolor=Salmon!5,linecolor=RawSienna,}
\declaretheoremstyle[headfont=\bfseries\color{RawSienna},
mdframed={style=mdredbox},headpunct={\\[3pt]},postheadspace=0pt,]{thmredbox}

\mdfdefinestyle{mdgreenbox}{linecolor=ForestGreen,backgroundcolor=ForestGreen!5,
	linewidth=2pt,rightline=false,leftline=true,topline=false,bottomline=false,}
\declaretheoremstyle[headfont=\bfseries\sffamily\color{ForestGreen!70!black},
mdframed={style=mdgreenbox},headpunct={ --- },]{thmgreenbox}

\mdfdefinestyle{mdorangebox}{linecolor=RedOrange,backgroundcolor=RedOrange!5,
	linewidth=2pt,rightline=false,leftline=true,topline=false,bottomline=false,}
\declaretheoremstyle[headfont=\bfseries\sffamily\color{RedOrange!70!black},
mdframed={style=mdorangebox},headpunct={ --- },]{thmorangebox}

\mdfdefinestyle{mdblackbox}{linecolor=black,backgroundcolor=RedViolet!5!gray!5,
	linewidth=3pt,rightline=false,leftline=true,topline=false,bottomline=false,}
\declaretheoremstyle[mdframed={style=mdblackbox}]{thmblackbox}

\declaretheoremstyle[spaceabove=6pt,spacebelow=6pt]{basehead}

\declaretheorem[style=basehead,name=Problem]{problem}
\declaretheorem[style=thmredbox,name=Problem,numbered=no]{problem*}
\declaretheorem[style=thmbloobox,name=Setup]{setup} % for config geo
\declaretheorem[style=thmredbox,name=Example,sibling=problem]{example}
\declaretheorem[style=thmredbox,name=$\bigstar$ Problem,sibling=problem]{niceprob}
\declaretheorem[style=thmbluebox,name=Theorem,sibling=problem]{theorem}
\declaretheorem[style=thmbluebox,name=Corollary,sibling=theorem]{corollary}
\declaretheorem[style=thmbluebox,name=Proposition,sibling=theorem]{prop}
\declaretheorem[style=thmbluebox,name=Lemma,numberwithin=problem]{lemma}
\declaretheorem[style=thmbluebox,name=Theorem,numbered=no]{theorem*}
\declaretheorem[style=thmbluebox,name=Lemma,numbered=no]{lemma*}
\declaretheorem[style=thmgreenbox,name=Claim,numberwithin=problem]{claim}
\declaretheorem[style=thmgreenbox,name=Claim,numbered=no]{claim*}
\declaretheorem[style=thmblackbox,name=Remark,numberwithin=problem]{remark}
\declaretheorem[style=thmblackbox,name=Remark,numbered=no]{remark*}
\declaretheorem[style=thmgreenbox,name=Definition,sibling=theorem]{definition}
\declaretheorem[style=thmgreenbox,name=Definition,numbered=no]{definition*}
\declaretheorem[style=thmorangebox,name=Warning,sibling=theorem]{warning}
\declaretheorem[style=thmorangebox,name=Warning,numbered=no]{warning*}
\declaretheorem[style=thmorangebox,name=Note,numbered=no]{note}
\declaretheorem[style=thmblackbox,name=Exercise,sibling=problem]{exercise}
\declaretheorem[style=thmblackbox,name=Exercise,numbered=no]{exercise*}
\declaretheorem[style=thmblackbox,name=Question,sibling=problem]{question}
\declaretheorem[style=thmblackbox,name=Question,numbered=no]{question*}

\addtolength{\textheight}{3.14cm}
\providecommand{\ol}{\overline}
\providecommand{\ul}{\underline}
\providecommand{\eps}{\varepsilon}
\providecommand{\half}{\frac{1}{2}}
\providecommand{\dang}{\measuredangle} %% Directed angle
\providecommand{\CC}{\mathbb C}
\providecommand{\FF}{\mathbb F}
\providecommand{\NN}{\mathbb N}
\providecommand{\QQ}{\mathbb Q}
\providecommand{\RR}{\mathbb R}
\providecommand{\ZZ}{\mathbb Z}
\providecommand{\dg}{^\circ}
\providecommand{\ii}{\item}
\providecommand{\setdiff}{\smallsetminus}
\providecommand{\alert}[1]{{\sffamily\textbf{\textcolor{blue}{#1}}}}
\providecommand{\inv}{^{-1}}
\providecommand{\mb}[1]{\mathbf{#1}}
\providecommand{\dif}{\;d}

\reversemarginpar

\newenvironment{soln}{\begin{proof}[Solution]}{\end{proof}}
\newenvironment{walkthrough}{\noindent \textbf{Walkthrough.}}{\medskip}
\newlist{walk}{enumerate}{3}
\setlist[walk]{label=\bfseries (\alph*)}

\providecommand{\pow}{\operatorname{Pow}}
\providecommand{\vocab}[1]{{\textbf{\textcolor{ForestGreen}{#1}}}}
\providecommand{\lcm}{\operatorname{lcm}}
\providecommand{\abs}[1]{\left|#1\right|}

\usepackage{asymptote}
\begin{asydef}
	defaultpen(fontsize(10pt));
	unitsize(1cm);
	size(8cm); // set a reasonable default
	usepackage("amsmath");
	usepackage("amssymb");
	settings.tex="pdflatex";
	settings.outformat="pdf";
	// Replacement for olympiad+cse5 which is not standard
	import geometry;
	// recalibrate fill and filldraw for conics
	void filldraw(picture pic = currentpicture, conic g, pen fillpen=defaultpen, pen drawpen=defaultpen)
	{ filldraw(pic, (path) g, fillpen, drawpen); }
	void fill(picture pic = currentpicture, conic g, pen p=defaultpen)
	{ filldraw(pic, (path) g, p); }
	// some geometry
	pair foot(pair P, pair A, pair B) { return foot(triangle(A,B,P).VC); }
	pair orthocenter(pair A, pair B, pair C) { return orthocentercenter(A,B,C); }
	pair centroid(pair A, pair B, pair C) { return (A+B+C)/3; }
	// cse5 abbreviations
	path CP(pair P, pair A) { return circle(P, abs(A-P)); }
	path CR(pair P, real r) { return circle(P, r); }
	pair IP(path p, path q) { return intersectionpoints(p,q)[0]; }
	pair OP(path p, path q) { return intersectionpoints(p,q)[1]; }
	path Line(pair A, pair B, real a=0.6, real b=a) { return (a*(A-B)+A)--(b*(B-A)+B); }
	// cse5 more useful functions
	picture CC() {
		picture p=rotate(0)*currentpicture;
		currentpicture.erase();
		return p;
	}
	pair MP(Label s, pair A, pair B = plain.S, pen p = defaultpen) {
		Label L = s;
		L.s = "$"+s.s+"$";
		label(L, A, B, p);
		return A;
	}
	pair Drawing(Label s = "", pair A, pair B = plain.S, pen p = defaultpen) {
		dot(MP(s, A, B, p), p);
		return A;
	}
	path Drawing(path g, pen p = defaultpen, arrowbar ar = None) {
		draw(g, p, ar);
		return g;
	}
\end{asydef}

\usepackage{mathrsfs}

\ihead{\footnotesize MAT 322 HW05}
\ohead{\footnotesize Last compiled \today}

\title{\vspace{-2em}MAT 322 HW05}
\author{\large Aditya Pahuja}
\date{\large Due 03/25}
\begin{document}
\maketitle

\begin{problem}
	\begin{enumerate}[(a)]
		\ii Suppose $f,g\colon Q\to\RR$ are Riemann integrable functions
		on a rectangle $Q\subseteq\RR^n$. Prove that
		$f+g$ is also Riemann integrable and that
		$\int_Qf+g=\int_Qf+\int_Qg$.
		\ii Show there exists a subset $S\subseteq\RR^2$ with measure zero
		whose closure and boundary both do not have measure zero.
		\ii Prove that non-empty open sets of $\RR^n$ do not have measure zero.
		\ii Let $f\colon[a,b]\to\RR$ be a continuous function.
		Prove that its graph
		\[G_f = \{(x,y)\in[a,b]\times\RR : y = f(x)\}\]
		is a subset of $\RR^2$ with measure zero.
	\end{enumerate}
\end{problem}

\begin{soln}
	(a) Given a partition $\mathcal P$ of $Q$, we see that
	\[U(f,\mathcal P)+U(g,\mathcal P)
	= \sum_{R\in\mathcal P} \sup_{x\in R}f(x) + \sup_{x\in R}g(x)
	\ge \sum_{R\in\mathcal P} \sup_{x\in R}(f(x)+g(x)) = U(f+g,\mathcal P)\]
	where the sums are over all rectangles in the partition; similarly
	we have the reverse inequality
	\[L(f,\mathcal P)+L(g,\mathcal P)\le L(f+g,\mathcal P)\]
	for the lower Riemann sums.
	Thus we have that for any partition $\mathcal P$ of $Q$,
	\[L(f,\mathcal P)+L(g,\mathcal P)
	\le L(f+g,\mathcal P)\le U(f+g,\mathcal P)
	\le U(f,\mathcal P)+U(g,\mathcal P).\]
	If $U(f,\mathcal P)-L(f,\mathcal P)$ and $U(g,\mathcal P)-L(g,\mathcal P)$
	are both less than $\eps/2$ (which we can arrange with a suitable
	choice of $\mathcal P$ due to Riemann integrability),
	then $U(f+g,\mathcal P)-L(f+g,\mathcal P)<\eps$,
	implying $f+g$ is Riemann integrable over $Q$.
	Furthermore, given $\eps>0$, if $\mathcal P$ is chosen so that
	\[U(f,\mathcal P)-\frac\eps2 < \int_Qf < L(f,\mathcal P)+\frac\eps2,\quad
	U(g,\mathcal P)-\frac\eps2 < \int_Qg < L(g,\mathcal P)+\frac\eps2,\]
	then the upper and lower Riemann sums of $f+g$ on $\mathcal P$
	are bounded between $\int_Qf+\int_Qg-\eps$ and $\int_Qf+\int_Qg+\eps$,
	so that
	\[\int_Qf+\int_Qg-\eps\le \int_Qf+g\le\int_Qf+\int_Qg+\eps\]
	for all $\eps>0$ implies that
	\[\int_Qf+\int_Qg = \int_Qf+g\]
	as desired.
	
	(b) Take $S = \QQ^2$. Since $S$ is countable, it has measure zero,
	and its interior is empty because any open ball contains points
	with all coordinates irrational, so the closure and boundary of $S$
	are both all of $\RR^2$. Clearly $\RR^2$ does not have measure zero.
	
	(c) We first show that rectangles $Q$ with nonzero volume
	do not have measure zero. Let $Q_1$, $Q_2$, \dots be
	a countable collection of rectangles whose union contains $Q$.
	Then, let $\eps > 0$ be any real number. Grow $Q_n$ into
	an open rectangle $R_n\supseteq Q_n$ whose volume exceeds $Q_n$
	by at most $\eps/2^n$, so that
	\[\sum_{n=1}^\infty\operatorname{vol}(R_n)
	\le \sum_{n=1}^\infty \operatorname{vol}(Q_n) + \frac{\eps}{2^n}
	= \eps + \sum_{n=1}^\infty \operatorname{vol}(Q_n).\]
	Now the $R_n$ form an open cover of the compact set $Q$,
	so some finite subcollection of the $R_n$ covers $Q$,
	say $R_{i_1}$, \dots, $R_{i_k}$. It is easy to show that
	\[\operatorname{vol}(Q)\le
	\operatorname{vol}\left(\bigcup_{j=1}^k R_{i_j}\right)
	\le \sum_{j=1}^k \operatorname{vol}(R_{i_k}),\]
	so we can conclude that
	\[\operatorname{vol}(Q)\le \eps+\sum_{n=1}^\infty\operatorname{vol}(Q_n)\]
	for all $\eps>0$, which implies that the total volume of the $Q_n$
	is at least as large as that of $Q$, hence the total volume of the $Q_n$
	is bounded below by a positive real number, which means $Q$
	cannot have measure zero.
	
	Now, let $U\subseteq\RR^n$ be open. We can find an open ball
	$B\subseteq U$, which in turn contains a rectangle $Q$ of nonzero volume.
	Any covering of $U$ with countably many rectangles $Q_n$
	must also cover $Q$, at which point the work above shows that
	the total volume of the $Q_n$ is bounded below by a positive real number,
	so $U$ also cannot have measure zero.
	
	(d) Since $[a,b]$ is compact, $f$ must be uniformly continuous.
	Given $\eps>0$, pick $\delta$ so that $|x-y|<\delta$
	for $x,y\in[a,b]$ implies $|f(x)-f(y)|<\frac{\eps}{b-a}$.
	Choose positive integer $n$ so that $\frac{b-a}n<\delta$.
	Then, on each interval $[a+\frac{k(b-a)}{n},a+\frac{(k+1)(b-a)}{n}]$
	for integer $0\le k<n$, the maximum and minimum of $f$
	on the interval differ by less than $\frac{\eps}{b-a}$,
	so that $G_f\cap([a+\frac{k(b-a)}{n},a+\frac{(k+1)(b-a)}{n}]\times\RR)$
	can be covered with a rectangle with dimensions
	$\frac{b-a}{n}\times\frac{\eps}{b-a}$ and volume $\frac{\eps}{n}$.
	This implies that $G_f$ can be covered by $n$ such rectangles
	for a total volume of $\eps$, hence $G_f$ can be covered by
	an at most countable collection of rectangles with volume at most $\eps$
	for any $\eps>0$, which is what it means for $G_f$ to have measure zero.
\end{soln}

\begin{problem}
	\begin{enumerate}[(a)]
		\ii Show that if $S\subseteq\RR^n$ is a compact subset of measure zero,
		then, for any $\eps > 0$, $S$ can be covered with
		finitely many rectangles with total volume less than $\eps$.
		\ii Let $Q\subseteq\RR^n$ be a rectangle and let $f\colon Q\to\RR$
		be a bounded function. Show that if $f$ vanishes outside
		a closed set of measure zero, then $\int_Qf$ exists and equals zero.
	\end{enumerate}
\end{problem}

\begin{soln}
	(a) By definition of $S$ having measure zero, there exists a
	countable collection $Q_1$, $Q_2$, \dots of rectangles
	whose union contains $S$ and whose total volume is less than $\eps/2$.
	For each $n$, let $R_n$ be an open rectangle containing $Q_n$
	whose volume is at most $Q_n + \eps/2^{n+1}$, so that
	the total volume of the $R_n$ exceeds that of $Q_n$ by at most $\eps/2$,
	hence is less than $\eps$. Then the $R_n$ form an open cover of $S$,
	so we can find a finite collection $R_{i_1}$, \dots, $R_{i_k}$
	such that their union contains $S$ and their total volume is at most
	$\sum_{n=1}^\infty \operatorname{vol}(R_n) < \eps$.
	This suffices; if we want closed rectangles then
	we just replace each $R_{i_j}$ with its closure, as a rectangle
	has the same volume as its interior.
	
	(b) Since $Z$ is closed and bounded,
	it is compact, so for a given $\eps>0$,
	it can be covered with finitely many boxes
	(which we will restrict to be contained in $Q$)
	whose total volume is at most $\eps$, by part (a).
	Let $\mathcal P$ be the partition induced by these boxes
	(so $\mathcal P$ is the coarsest partition such that each of these boxes
	are unions of some rectangles in $\mathcal P$).
	If $m \le f \le M$ are lower and upper bounds for $f$ on $Q$, then
	\[m\eps \le L(f,\mathcal P) \le U(f,\mathcal P)\le M\eps\]
	since the total volume of the rectangles in $\mathcal P$ that intersect $Z$
	is at most $\eps$ by construction of $\mathcal P$,
	and all the other rectangles contribute zero to the Riemann sums.
	In particular
	\[m\eps\le \underline\int_Qf\le\ol\int_Qf\le M\eps
	\implies \underline\int_Qf=\ol\int_Qf=\int_Qf=0\]
	by taking $\eps\to 0$,
	so $f$ is Riemann integrable over $Q$ and has integral zero.
\end{soln}

\begin{problem}
	\begin{enumerate}[(a)]
		\ii Let $f,g\colon S\to\RR$ be bounded functions on a bounded set
		$S\subseteq\RR^n$. Let
		\[M(x) = \max(f(x),g(x)),\quad m(x)=\min(f(x),g(x)).\]
		Show that if $f$ and $g$ are continuous
		(resp. Riemann integrable over $S$)
		then so are $M$ and $m$.
		\ii Let $S = S_1\cup S_2\subseteq\RR^n$ be a bounded set,
		and $f\colon S\to\RR$ be a function that is
		Riemann integrable over $S_1$ and $S_2$. Prove that
		$f$ is Riemann integrable over $S$ and $S_1\cap S_2$ and
		\[\int_Sf=\int_{S_1}f+\int_{S_2}f-\int_{S_1\cap S_2}f.\]
	\end{enumerate}
\end{problem}

\begin{soln}
	(a) For continuity, we can write
	\[\max(f(x),g(x)) = \frac{f(x) + g(x) + |f(x)-g(x)|}{2},\]
	since if WLOG $f(x)\ge g(x)$ then
	\[\frac{f(x)+g(x)+|f(x)-g(x)|}{2}
	= \frac{f(x)+g(x)+f(x)-g(x)}{2} = f(x) = \max(f(x),g(x)).\]
	Linear combinations and compositions preserve continuity
	(noting that $x\mapsto |x|$ is continuous),
	so this is continuous because $f$ and $g$ are.
	Then since $\min(f(x),g(x))=f(x)+g(x)-\max(f(x),g(x))$,
	$m(x)$ is a linear combination of continuous functions, hence continuous.
	
	For Riemann integrability, we start by observing that
	if $f$ and $g$ are continuous at $x$, then so are $m$ and $M$.
	If $f$ and $g$ are Riemann integrable over $S$,
	$f$ and $g$ are only discontinuous on a set of measure zero
	within some rectangle $Q\supseteq S$ (where we extend $f$ and $g$
	to be zero outside $S$).
	The union of the sets of discontinuities of $f$ and $g$
	has measure zero because the union of two measure zero sets
	has measure zero, and the set of discontinuities of $M$ and $m$
	is a subset of this union, so $M$ and $m$ are also
	discontinuous on a set of measure zero, hence Riemann integrable over $S$.
	
	(b) Note that we can write the positive part of $f$ as
	$f_+=\frac{f+|f|}{2}$, so $f_+$ is Riemann integrable on $S_1$ and $S_2$.
	Similarly, so is the negative part $f_-=\frac{|f|-f}{2}$.
	Now, $f_+ = \max({f_+}_{S_1},{f_+}_{S_2})$, so by (a)
	$f_+$ is Riemann integrable on $S$ (since the restricted functions
	are also Riemann integrable on $S$ by virtue of being zero
	outside $S_1$ and $S_2$ respectively).
	Analogously, ${f_+}_{S_1\cap S_2} = \min({f_+}_S_1,{f_+}_S_2)$
	implies that $f_+$ is Riemann integrable on $S_1\cap S_2$.
	We can also conclude that $f_-$ is integrable on $S$ and $S_1\cap S_2$,
	hence the linear combination $f_+-f_- = f$ is integrable
	on $S$ and $S_1\cap S_2$.
	
	Furthermore, we can rewrite
	\[\int_Ef = \int_S f\cdot 1_E\]
	for any $E\subseteq S$ on which $f$ is Riemann integrable, so
	proving the identity of integrals amounts to showing that
	\[f = f\cdot 1_{S_1} + f\cdot 1_{S_2} - f\cdot 1_{S_1\cap S_2}\]
	by the linearity of the Riemann integral,
	and this identity is trivial, so the integral identity holds.
\end{soln}

\begin{problem}
	\begin{enumerate}[(a)]
		\ii Suppose $S_1,S_2\subseteq\RR^n$ are rectifiable sets.
		Prove that $S_1\cup S_2$ and $S_1\cap S_2$ are rectifiable,
		and their volumes are related by
		\[\operatorname{vol}(S_1) + \operatorname{vol}(S_2)
		= \operatorname{vol}(S_1\cup S_2)-\operatorname{vol}(S_1\cap S_2).\]
		\ii Suppose $S\subseteq\RR^n$ is a bounded set
		which is a union of countably many rectifiable sets
		$S_1$, $S_2$, $S_3$, \dots. Prove or disprove that
		$S$ must be rectifiable.
	\end{enumerate}
\end{problem}

\begin{soln}
	(a) Note that a bounded set $S$ is rectifiable
	if and only if there exists a rectangle $Q\supseteq S$
	such that the function $1_S\colon Q\to\RR$ defined by
	\[1_S(x) := \begin{cases}
		1&x\in S\\
		0&x\notin S
	\end{cases}\]
	is Riemann integrable over $S$, as the discontinuities of $1_S$
	are precisely the points on the boundary of $S$, and Riemann integrability
	is equivalent to this set of discontinuities having measure zero.
	Furthermore,
	\[\operatorname{vol}(S) = \int_S 1_S.\]
	Problem 3(b) now applies: we know that $1_{S_1\cup S_2}$
	is Riemann integrable over $S_1$ and $S_2$, so it must be
	Riemann integrable over $S_1\cup S_2$ and $S_1\cap S_2$ as well, with
	\begin{align*}
		\operatorname{vol}(S_1\cup S_2)
		&= \int_{S_1\cup S_2} 1_{S_1\cup S_2}\\
		&= \int_{S_1} 1_{S_1\cup S_2} + \int_{S_2} 1_{S_1\cup S_2}
		- \int_{S_1\cap S_2}1_{S_1\cup S_2}\\
		&= \operatorname{vol}(S_1) + \operatorname{vol}(S_2)
		- \operatorname{vol}(S_1\cap S_2),
	\end{align*}
	which rearranges to the desired expression for
	the sum of the volumes of $S_1$ and $S_2$.
	
	(b) $S$ need not be rectifiable. For example, consider an enumeration
	$(q_n)_{n\in\NN}$ of $\QQ\cap[0,1]$ and define
	$S_n = \{q_1,q_2,\dots,q_n\}$.
	Each $S_n$ is a finite set, so the boundary of $S_n$ is $S_n$
	which clearly has measure zero (given $\eps>0$,
	we can just place a rectangle of volume $\eps/n$ on each point of $S_n$);
	that is, the $S_n$ are all rectifiable.
	However, the union of the $S_n$ is $S = \QQ\cap[0,1]$,
	which isn't rectifiable because the closure of $S$ is $[0,1]$
	while the interior is empty, so that the boundary of $S$ is $[0,1]$,
	which clearly doesn't have measure zero.
\end{soln}

\begin{problem}
	In this question we only consider integrals in the ``extended sense.''
	\begin{enumerate}[(a)]
		\ii Prove that $\int_{\RR^2}e^{-x^2-y^2}$ exists and equals $\pi$.
		\ii Prove that the \vocab{Gaussian integral} $\int_\RR e^{-x^2}$
		exists and equals $\sqrt\pi$.
	\end{enumerate}
\end{problem}


\begin{soln}
	(a) Consider the increasing sequence
	$K_1\subseteq K_2\subseteq\cdots\subseteq\RR^2$ of compact rectifiable sets
	where $K_n$ is the circle of radius $n$ centered at $(0,0)$.
	The union of the $K_n$ is all of $\RR^2$. Furthermore,
	by a change of coordinates, we can compute
	\[\int_{K_n}e^{-x^2-y^2} = \int_{(0,2\pi)\times(0,n)} e^{-r^2}\cdot r,\]
	where $r$ (which varies across $(0,n)$
	is the determinant of the Jacobian of the diffeomorphism
	that changes Cartesian coordinates to polar coordinates
	(actually the RHS is the integral over $K_n$ minus
	its boundary and the positive $x$-axis, but these sets
	contribute zero to the integral anyway).
	The inner integral exists independently of the first coordinate
	because it is continuous in $r$, so Fubini's theorem gives
	\[\int_{(0,2\pi)\times(0,n)} e^{-r^2}\cdot r =
	\int_{\theta\in (0,2\pi)}\left(\int_{r\in(0,n)} re^{-r^2}\right)
	= \pi\cdot \left(1 - e^{-n^2}\right).\]
	This means
	\[\int_{\RR^2} e^{-x^2-y^2} =
	\lim_{n\to\infty}\int_{K_n} e^{-x^2-y^2} = \pi < \infty,\]
	so the integral over $\RR^2$ exists and equals $\pi$.
	
	(b) This integral exists because, noting that $e^{-x^2}$
	is continuous on $\RR$, it is Riemann integrable on $[-n,n]$ for all $n$.
	Furthermore, we can observe the inequality
	\[e^{-m^2} \le e^{-m} \le 2^{-n}\]
	for all nonnegative integers $m$.
	If $m(x)$ is the floor of $|x|$, this lets us bound $e^{-x^2}\le 2^{-m(x)}$.
	We can calculate for all positive integers $n$
	\[\int_{[-n,n]} 2^{-m(x)} = 2\sum_{k=0}^{n-1} 2^{-k} = 2(2 - 2^{-n+1})\]
	(integrability follows from being piecewise constant, so $2^{-m(x)}$
	has finitely many discontinuities)
	so taking the limit of $n$ shows that $2^{-m(x)}$ is integrable on $\RR$
	with integral equal to $4$. Thus
	\[\lim_{n\to\infty} \int_{[-n,n]}e^{-x^2}
	\le \lim_{n\to\infty} \int_{[-n,n]}2^{-m(x)} = 4 < \infty\]
	implies that $e^{-x^2}$ is integrable over $\RR$.
	Now we can use Fubini's theorem to get
	\[\int_{[-n,n]\times[-n,n]} e^{-x^2-y^2}
	= \int_{[-n,n]}e^{-x^2}\cdot \int_{[-n,n]}e^{-y^2}
	= \left(\int_{[-n,n]}e^{-x^2}\right)^2,\]
	so that taking $n\to\infty$ gives
	\[\pi = \int_{\RR^2}e^{-x^2-y^2} = \left(\int_\RR e^{-x^2}\right)^2,\]
	which shows that the Gaussian integral is $\sqrt\pi$.
\end{soln}
























\end{document}